\setcounter{subsection}{0}
\section{DISEÑO DE LA SOLUCIÓN}
%Debe presentar la arquitectura de la solución y los principales artefactos del diseño detallado, junto con la descripción de las herramientas y tecnologías seleccionadas para resolver el pro-blema. Se recomienda justificar la selección de herramientas basándose en múltiples criterios (establecidos durante la fase de análisis) mediante el uso de tablas comparativas.
%Para más detalles, recomendamos consultar documentos adjuntos que expliquen todos los as-pectos del diseño (por ejemplo, SDD).
%Recuerde que diseñar es un proceso previo a la implementación.

\subsection{Protocolo de evaluación de M3}
M3 evalúa los candidatos que entrega M2 sobre la reserva: el 20\,\% final y cronológico del panel, posterior a la fecha de corte $t^*$. Cada candidato se ajusta una sola vez con la historia hasta $t^*$ y no se reajusta. Desde cada origen diario de la reserva, cada modelo pronostica los siete días siguientes, y esos siete pronósticos forman una ventana.

\subsubsection{Bloques de selección y de prueba}
M3 divide la reserva en dos bloques consecutivos. El bloque de selección, el 40\,\% inicial de los días reservados, elige la familia campeona de cada clase de demanda y la congela. El bloque de prueba, el 60\,\% final, no interviene en ninguna decisión y produce todos los resultados reportados. Los seis orígenes cuyas ventanas cruzan el límite entre los bloques no se usan. Así, el campeón nunca se califica con los datos que lo eligieron, una fuente conocida de optimismo en la evaluación \cite{cawley2010}. Las fracciones se fijaron antes de las corridas finales y no se ajustaron con sus resultados.

\subsubsection{Productos que compiten}
\begin{sloppypar}
Un producto entra en la competencia solo si al menos dos familias producen un modelo válido en M2. Un producto con una sola familia evaluable sale de la competencia con la causa \texttt{eligibility.\allowbreak fewer\_than\_two\_families}. Un producto sin ninguna familia evaluable recibe el veredicto \texttt{FALLO\_TECNICO} y se cuenta en su clase. Estas reglas dependen solo de resultados anteriores a la reserva.
\end{sloppypar}

\subsubsection{Selección y confirmación del campeón}
En el bloque de selección, M3 elige por clase de demanda la familia con menor mediana de la razón $r = RMSE_{\text{modelo}} / RMSE_{\text{Seasonal Naive}}$, calculada sobre el mismo conjunto de productos para todas las familias. Si esa mediana no es menor que 1, el campeón es Seasonal Naive. Si dos familias difieren en menos de una tolerancia predeclarada, gana la más simple.

En el bloque de prueba, M3 confirma el campeón congelado frente a cada familia competidora. Para cada producto calcula el logaritmo de la razón entre el RMSE del campeón y el del competidor, y aplica la prueba de rangos con signo de Wilcoxon \cite{wilcoxon1945} sobre esos valores. El campeón debe superar a todos los competidores de su clase, y la corrección de Holm \cite{holm1979} ajusta los valores p de las cuatro clases. Además, un intervalo bootstrap por bloques de 14 días de la mediana de la razón debe quedar por debajo de 1. Una confirmación fallida se reporta, pero no cambia el campeón. La interpretación de Wilcoxon como una comparación de localización supone independencia entre productos y simetría de los logaritmos de las razones bajo la hipótesis nula. El bootstrap conserva los choques comunes al remuestrear los mismos días para todos los productos. Sus intervalos suponen que el bloque de 14 días representa la dependencia temporal relevante.

\subsubsection{Prueba frente a Seasonal Naive y veredictos}
En el bloque de prueba, M3 compara el campeón de cada producto con Seasonal Naive mediante la prueba HLN-DM \cite{dm,harvey1997}. Su interpretación supone un diferencial de pérdida estacionario y dependencia despreciable más allá del rezago seis. La corrección de Benjamini-Hochberg \cite{benjamini1995} se aplica dentro de cada clase de demanda con nivel nominal $q = 0{,}05$. Su control de la tasa de descubrimientos falsos requiere independencia o la condición de dependencia positiva de Benjamini-Yekutieli. Los choques compartidos entre productos no bastan para establecer esa condición. La corrección de Benjamini-Yekutieli se reporta como sensibilidad válida bajo dependencia arbitraria entre valores p válidos.

Cada producto recibe uno de cinco veredictos descriptivos, calculados en el bloque de prueba: \texttt{FALLO\_TECNICO}, \texttt{NO\_EVALUABLE}, \texttt{EVIDENCIA\_INSUFICIENTE}, \texttt{VALIDADO} o \texttt{EXPLORATORIO}. \texttt{VALIDADO} exige datos reales, al menos 30 ventanas válidas y $r < 1$, pero no exige significancia. La significancia se muestra en su propia columna. Por eso, las tablas de resultados presentan dos conteos por clase: los productos \texttt{VALIDADO} y los productos con una mejora significativa frente a Seasonal Naive. Una diferencia no significativa no se interpreta como equivalencia.

\subsubsection{Clases de demanda y pronóstico puntual de Chronos-2}
La clase Syntetos-Boylan de cada producto se calcula con toda su historia hasta $t^*$ y no cambia en la reserva. Un producto puede comportarse en la reserva de forma distinta a su clase. Por eso, los resultados de cada clase se presentan junto con el perfil de demanda que sus productos muestran en el bloque de prueba.

Chronos-2 emite cuantiles, y PRED toma la mediana como su pronóstico puntual. Con pérdida cuadrática, la media de la distribución condicional verdadera minimiza el error esperado \cite{gneiting2011}. Esta propiedad no garantiza que la media estimada por un modelo supere a su mediana en los datos evaluados. PRED mantiene la mediana como elección predeclarada y publica el sesgo con signo de cada modelo por producto.

\subsubsection{Grupo final de corridas}
Los resultados de la tesis provienen de un grupo de tres corridas con las semillas 0, 1 y 2, fijadas en el protocolo antes de la primera corrida. Las corridas difieren solo en la semilla del motor. Comparten el código, los datos, la lista de productos, los bloques de la reserva y la política. Las corridas cubren el catálogo completo de productos del panel si el control de memoria de la etapa de ajuste lo permite. Si no lo permite, cubren la muestra fija más grande que ese control admite, la misma en las tres corridas. La semilla cambia la búsqueda de hiperparámetros de las familias optimizadas, así que el grupo mide si la clasificación de las familias se mantiene ante esa aleatoriedad. Un campeón es estable si gana su clase y pasa su confirmación en las tres corridas. La significancia del grupo usa, para cada producto que compite y tiene una prueba válida en las tres corridas, el mayor de sus tres valores p bilaterales, con una sola corrección de Benjamini-Hochberg por clase. Los productos sin alguna prueba válida se reportan con su causa. Una mejora significativa del grupo exige además un diferencial de pérdida negativo frente a Seasonal Naive en las tres corridas. Las corridas de desarrollo no forman parte de la evidencia.
